% !TEX root = thesis.tex

%% Document Metadata

%% The location of your Assignment Statement (scanned as image)
\thesisspec{figures/zadavaci-list.png}

\title{Quantum Random Number Generator}{Generovanie náhodných čísel pomocou kvantového počítača}
\thesis{Bachelors thesis}{Bakalárska práca}  % typ prace

\author{Peter}{Jaš}
\supervisor{Ing. Miroslav Murin}   % thesis supervisor
% \consultant{Donald E. Knuth}  % thesis consultant
% \consultantsaffiliation{Stanford University}{Stanfordská univerzita}  % affiliation of the consultant
 \college{Technical University of Košice}{Technická univerzita v Košiciach}  % univerzita
% \faculty{Faculty of Electrical Engineering and informatics}{Fakulta elektrotechniky a informatiky}  % fakulta
 \department{Department of Computers and Informatics}{Katedra počítačov a informatiky}  % katedra
 \departmentacr{DCI}{KPI}  % skratka katedry
\submissiondate{29}{5}{2026}
 \fieldofstudy{Informatika}
 \studyprogramme{Informatika}  % iba názov bez číselného kódu
% \city{}{Košice}  % mesto

\keywords{%
% english
  quantum computing, random number generation, quantum randomness, Bell pairs,
  qubit entanglement, QRNG, entropy extraction, next-bit test
}{%
% slovak
  kvantové počítače, generovanie náhodných čísel, kvantová náhodnosť, Bellove páry,
  prepletenie qubitov, QRNG, extrakcia entropie, test nasledujúceho bitu
}

 \declaration{
    Vyhlasujem, že som záverečnú prácu vypracoval(a) samostatne s použitím uvedenej odbornej literatúry, ako aj s využitím nástrojov generatívnej umelej inteligencie v súlade s Metodickým pokynom o záverečných a kvalifikačných prácach a Etickým kódexom študenta Technickej univerzity v Košiciach. Umelá inteligencia bola použitá ako podporný nástroj bez porušenia zásad akademickej etiky a autorskej integrity.
 }

\abstract{%
% english
  This thesis presents a quantum random number generator (Q.R.N.G.) based on Bell
  pair measurements. Entangled qubit pairs prepared through Hadamard and CNOT gates produce
  correlated measurement outcomes used as a source of random bits, which are
  stored in a database and exposed through a Django web application. Output quality
  was evaluated using NIST-inspired tests, next-bit prediction, and Markov
  analysis — the generator passed all tests with no detectable deterministic
  structure. The thesis also examines the
  impact of hardware imperfections such as decoherence and measurement bias, and
  demonstrates that entropy extraction via MD5 hashing reduces local imbalances
  without introducing artificial randomness.
}{%
% slovak
  Táto práca sa zaoberá generovaním náhodných čísel pomocou kvantového počítača
  so zameraním na praktickú využiteľnosť kvantového prepletenia ako zdroja
  skutočnej náhodnosti. Bol navrhnutý a implementovaný kvantový generátor
  náhodných čísel založený na meraniach Bellových párov. Generátor využíva dvojice
  prepletených qubitov pripravených pomocou Hadamardovej a CNOT brány, pričom
  výsledné bity sú ukladané do databázy a sprístupnené cez webovú aplikáciu
  postavenú na frameworku Django. Kvalita výstupu bola overená testami
  inšpirovanými metodikou NIST, testom predikcie nasledujúceho bitu a Markovovým
  testom. Generátor úspešne prešiel všetkými testami bez detekovateľnej
  deterministickej štruktúry. Práca sa venuje aj vplyvu nedokonalostí
  kvantového hardvéru, ako sú decoherencia a bias merania, a preukazuje, že
  extrakcia entropie pomocou hašovacej funkcie MD5 redukuje lokálne nerovnováhy
  bez umelého zvyšovania entropie.
}

\acknowledgment{
  Na tomto mieste by som sa rád poďakoval svojmu vedúcemu práce, Ing. Miroslavovi Murinovi, za jeho čas a odborné vedenie počas riešenia mojej záverečnej práce.
  Rovnako by som sa rád poďakoval svojim rodičom a priateľom za ich podporu a povzbudzovanie počas celého môjho štúdia.
  Osobitná vďaka patrí Jakubovi Novákovi za inšpiratívny nápad a prof. Leonardovi Francovi z Univerzity v Málage za cenné rady ohľadom systematického hodnotenia náhodnosti a štruktúrovania vedeckej práce.
}


% Uncomment the following line for the printer-friendly version of the thesis.
% When defined, final PDF will not contain colored links, which is useful for printing.
% WARNING: Uncomment ONLY when used manually without the Docker image for building the thesis. By uncommenting this line, you will manually override the printable flag and behavior of the Docker image will not work as expected.
% If you want to use the Docker image for building the document, please do not uncomment this line.
%\def\printable{}

